Since this PCB housed several circuits, it was important to use the silk-screen to promote good labelling. Circuits like microSD connector were enclosed in a chamfered rectangle to show which components contribute to the circuit. All test points used the 'Arial Black' font instead of the 'Arial Narrow' so I can easily identify where the testpoints are during manufacturing and testing. You may find jokes on the silkscreen like "BRAIN GOES HERE", these hold no importance towards the actual PCB itself, they have just been included for entertainment purposes
Once the components were placed, these jokes were hidden. Silkscreen preparation was the last step of the PCB, so if you explore previous saves of this PCB, you will most likely see lots of silkscreen collisions.
\nl
To save space on the PCB, some test point labels were shortened. Additionally, the test points inside the power LED array (i.e., next to 3V3, 2V8 and 1V8) are the drain testpoints (there was not enough space to include the silkscreen).

\begin{table}
    \centering
    \begin{tabular}{|l|l|}
        \hline
        \textbf{PCB Name} & \textbf{Full Name} \\
        \hline
        SCRL\_BTN & Scroll Button \\ \hline
        R3V3 & Regulated 3V3 \\ \hline
        PA & Power Action \\ \hline
        PS & Power State \\ \hline
        DBG & Debug \\ \hline
        DRN & Drain \\ \hline
        SLED & Sensor LED \\ \hline
        MLED & MCU LED \\ \hline
        RLED & Read LED \\ \hline
        WLED & Write LED \\ \hline
        VTMH & Vibration Table Mounting Hole \\ \hline
    \end{tabular}
    \caption{Explanation of truncated silkscreen labels}
\end{table}